\documentclass[10pt,a4paper]{article}
\usepackage{index-flmx.37ca24e7}
\usepackage[hmargin=21mm,top=19mm,bottom=20mm,includeheadfoot,headheight=24pt,headsep=7mm,footskip=10mm,marginparwidth=15mm,marginparsep=3mm]{geometry}
\usepackage{fontspec}
% Load the font files and weight/italic mappings supplied by TeX Live.
\usepackage[oldstyle,proportional]{sourceserifpro}
\usepackage[semibold]{sourcesanspro}
\usepackage[lining,regular,bold]{sourcecodepro}
\usepackage{unicode-math}
\setmathfont{STIX Two Math}
% Supply two characters emitted by FLM that Source Serif Pro can lack:
% the narrow nonbreaking space in references and the proof's closing square.
\usepackage{newunicodechar}
\newunicodechar{ }{\nobreak\thinspace}
\newunicodechar{□}{\ensuremath{\mdlgwhtsquare}}
% Hanging punctuation and restrained font expansion improve justification.
% Expansion limits are in thousandths (15 = 1.5%); keep monospaced text out.
\usepackage[protrusion=true,expansion=true,stretch=15,shrink=15,step=1]{microtype}
\UseMicrotypeSet[expansion]{alltext-nott}
\usepackage{fancyhdr}
\usepackage{truncate}
\usepackage{framed}

\definecolor{memoAccent}{HTML}{245C63}
\definecolor{memoInk}{HTML}{252D30}
\definecolor{memoMuted}{HTML}{657477}
\colorlet{memoRule}{memoAccent!25!white}
\colorlet{shadecolor}{memoAccent!5!white}
\colorlet{flmxLinkColor}{memoAccent}
\colorlet{flmxUrlColor}{memoAccent}
\colorlet{flmxCiteColor}{memoAccent}

% Smaller text with more leading relative to its size, and modest paragraph gaps.
\linespread{1.18}
\setlength{\parindent}{0pt}
\setlength{\parskip}{0.48em plus 0.12em minus 0.08em}
\setlength{\emergencystretch}{1.5em}
\setlength{\textfloatsep}{14pt plus 3pt minus 2pt}
\setlength{\intextsep}{12pt plus 3pt minus 2pt}
\setlength{\abovecaptionskip}{6pt}
\setlength{\FrameSep}{9pt}
\raggedbottom
\clubpenalty=3000
\widowpenalty=3000

% Section marks are provided by FLM, including alternative heading text.
% Abbreviate unusually long running heads to keep the header on one line.
\pagestyle{fancy}
\fancyhf{}
%\fancyhead[L]{\sffamily\fontsize{8}{10}\selectfont\color{memoMuted}\truncate{\textwidth}{\nouppercase{\rightmark}}}
%\fancyfoot[R]{\sffamily\fontsize{9}{11}\selectfont\color{memoAccent}\thepage}
\fancyfoot[C]{\sffamily\fontsize{9}{11}\selectfont\color{memoAccent}\thepage}
\renewcommand{\headrulewidth}{0pt}%.3pt}
\renewcommand{\footrulewidth}{0pt}
%\renewcommand{\headrule}{\hbox to\headwidth{\color{memoRule}\leaders\hrule height\headrulewidth\hfill}}
\fancypagestyle{memoFirst}{\fancyhead{}\renewcommand{\headrulewidth}{0pt}}

% Keep anchors in every block whose default look is replaced.
\newcommand{\memoAnchor}{\flmxIfKey{anchor}{\flmxTarget{\flmxKey{anchor}}}{}}

% Metadata stays in FLM's rendered form, so emphasis, mathematics and links
% in titles and author names retain their original meaning.
\newcommand{\memoMetadataFont}{\sffamily\mdseries\upshape\fontsize{9.5}{12}\selectfont\color{memoInk}}
\flmxDefineBlock{flm.meta.title}
  {\par\memoAnchor\raggedright\sffamily\bfseries\fontsize{20}{25}\selectfont\color{memoInk}}
  {\par\vspace{10pt}}
\flmxDefineBlock{flm.meta.author}
  {\par\memoAnchor\raggedright\memoMetadataFont
    \leftskip=2em\relax}
  {\par\vspace{5pt}}
\flmxDefineBlock{flm.meta.date}
  {\par\memoAnchor\raggedright\memoMetadataFont
    \leftskip=2em\relax}
  {\par}

% Delegate placement to FLM: numbering, targets, PDF bookmarks and running
% heads remain intact. Change only the type within each heading level.
\flmxDefineHeading{0}{\flmxHeadingLevel{0}{{\sffamily\fontsize{19}{23}\selectfont\color{memoAccent}#1}}}
\flmxDefineHeading{1}{\flmxHeadingLevel{1}{{\sffamily\fontsize{14}{18}\selectfont\color{memoAccent}#1}}}
\flmxDefineHeading{2}{\flmxHeadingLevel{2}{{\sffamily\fontsize{11.5}{15}\selectfont\color{memoInk}#1}}}
\flmxDefineHeading{3}{\flmxHeadingLevel{3}{{\sffamily\normalsize\color{memoInk}#1}}}
\flmxDefineHeading{4}{\flmxHeadingLevel{4}{{\sffamily\normalsize\color{memoInk}#1}}}
\flmxDefineHeading{5}{\flmxHeadingLevel{5}{{\sffamily\normalsize\color{memoInk}#1}}}
\flmxDefineHeading{xnotes}{\flmxHeadingLevel{2}{{\sffamily\small\color{memoMuted}#1}}}
\flmxDefineHeading{citations}{\flmxHeadingLevel{2}{{\sffamily\small\color{memoMuted}#1}}}

\makeatletter
\def\flmxHeadingLevel#1#2{%
  \ifcase#1\relax % case #1 = 0
    \par\addpenalty{-\@highpenalty}%
    \@startsection{flmxheading}{0}{\z@}%
      {-4.5ex \@plus -2ex \@minus -.4ex}{1.8ex \@plus .2ex}%
      {\normalfont\huge\bfseries}*{\flmx@headingplace #2}%
  \or % case #1 = 1
    \@startsection{flmxheading}{1}{\z@}%
      {-3.5ex \@plus -2ex \@minus -.4ex}{1.8ex \@plus .2ex}%
      {\normalfont\Large\bfseries}*{\flmx@headingplace #2}%
  \or % case #1 = 2
    \@startsection{flmxheading}{2}{\z@}%
      {-3.25ex \@plus -1ex \@minus -.2ex}{1.3ex \@plus .2ex}%
      {\normalfont\large\bfseries}*{\flmx@headingplace #2}%
  \or % case #1 = 3
    \@startsection{flmxheading}{3}{\z@}%
      {-3.25ex \@plus -1ex \@minus -.2ex}{1.1ex \@plus .2ex}%
      {\normalfont\normalsize\bfseries}*{\flmx@headingplace #2}%
  \or % case #1 = 4
    \@startsection{flmxheading}{4}{\z@}%
      {3.0ex \@plus .8ex \@minus .2ex}{-1em}%
      {\normalfont\normalsize\bfseries}*{\flmx@headingplace #2}%
  \or % case #1 = 5
    \@startsection{flmxheading}{5}{\parindent}%
      {3.25ex \@plus 1ex \@minus .2ex}{-1em}%
      {\normalfont\normalsize\bfseries}*{\flmx@headingplace #2}%
  \else
    [UNKNOWN HEADING LEVEL]
  \fi
}
\makeatother


% End-note and bibliography entries, including displays, use the smaller size.
\flmxDefineBlock{xnotes}
  {\par\memoAnchor\small}
  {\par}
\flmxDefineBlock{flm.citations.bibliography}
  {\par\memoAnchor\small}
  {\par}

% Plain, symmetric quotation indents; native lists can break across pages.
\flmxDefineBlock{flm.quote}
  {\begin{list}{}{\setlength{\leftmargin}{1.75em}\setlength{\rightmargin}{1.75em}%
     \setlength{\listparindent}{0pt}\setlength{\parsep}{\parskip}%
     \setlength{\topsep}{0.4\baselineskip}}%
   \item\relax\memoAnchor\color{memoMuted}}
  {\end{list}}
\flmxDefineBlock{flm.quote.attributed}
  {\par\memoAnchor\raggedleft\sffamily\footnotesize\color{memoMuted}}
  {\par}
\flmxDefineBlock{flm.floats.float.caption}
  {\par\vspace{\abovecaptionskip}\memoAnchor
    \sffamily\small\color{memoMuted}}
  {\par}

% Only override verbatim hooks when the generated package provides them:
% a document without code does not necessarily load this FLM component.
\ifdefined\flmxVerbatimFont
  \renewcommand{\flmxVerbatimFont}{\ttfamily\small}
  \renewcommand{\flmxVerbatimBlockBegin}{\begin{snugshade*}}
  \renewcommand{\flmxVerbatimBlockEnd}{\end{snugshade*}}
\fi

\begin{document}
\color{memoInk}
\thispagestyle{memoFirst}
% The first page has no running head; give its space back to the title.
\vspace*{\dimexpr-\headheight-\headsep\relax}
\noindent\begin{minipage}{\linewidth}
\setlength{\parskip}{0pt}
{\color{memoAccent}\hrule height 1.0pt\relax}
\vspace{12pt}
\begin{flmxBlock}[role=flm.meta.title, kinds={meta}]
A document of several files
\end{flmxBlock}

\begin{flmxBlock}[role=flm.meta.author, kinds={meta}]
\flmxSpan[role=flm.meta.author.name]{Ann Author}
\end{flmxBlock}

\begin{flmxBlock}[role=flm.meta.date, kinds={meta}]
30 September 2026
\end{flmxBlock}

\vspace{10pt}
{\color{memoAccent}\hrule height 0.8pt\relax}
\end{minipage}\par\vspace{4pt}
\flmxHeading[kind=section, level=1, toc=true, anchor=sec-intro, plain={Introduction}]{Introduction}

This paragraph begins in the primary file, continues with words of an included file, which names the \flmxSpan[style=emph]{Example}
project through the house package’s macro, and ends in the primary file, as one paragraph. The method
is in \flmxLink[kind=ref, href={\#sec-method}, target=sec-method]{The \flmxSpan[style=emph]{Method}}.\flmxSpan[style=superscript, role=flm.xnotes.xnote.mark, anchor=fnref--1yuhk]{\flmxLink[kind=xnote, href={\#fn-1}, target=fn-1, doc-role=noteref, label={Footnote a}]{a}}

\flmxHeading[kind=section, level=1, toc=true, anchor=sec-method, plain={The Method}]{The \flmxSpan[style=emph]{Method}}

A section kept in a file of its own, back to \flmxLink[kind=ref, href={\#sec-intro}, target=sec-intro]{Introduction}.\flmxSpan[style=superscript, role=flm.xnotes.xnote.mark, anchor=fnref--2yurg]{\flmxLink[kind=xnote, href={\#fn-2}, target=fn-2, doc-role=noteref, label={Footnote b}]{b}}

\flmxHeading[kind=subsection, level=2, toc=true, anchor=sec--Details, plain={Details}]{Details}

The details, read relative to the file that names them:

\begin{flmxList}[kinds={unnumbered itemize labeled}, depth=1]
\flmxItem[label={•}, anchor=item--1otnn]
one;

\flmxItem[label={•}, anchor=item--2otxj]
two.
\end{flmxList}
\flmxHeading[kind=section, level=1, toc=true, anchor=sec--Conclusion, plain={Conclusion}]{Conclusion}

The primary file goes on after the section kept in a file of its own, and its notes are counted after
that file’s.\flmxSpan[style=superscript, role=flm.xnotes.xnote.mark, anchor=fnref--3yvbc]{\flmxLink[kind=xnote, href={\#fn-3}, target=fn-3, doc-role=noteref, label={Footnote c}]{c}}

\begin{flmxBlock}[role=flm.xnotes.footnote.list, kinds={xnotes footnote}, doc-role=endnotes]
\flmxHeading[kind=xnotes, level=1, toc=false, plain={Footnotes}]{Footnotes}
\begin{flmxList}[kinds={numbered xnotes footnote labeled}, depth=1]
\flmxItem[label={\flmxSpan[style=superscript, role=flm.xnotes.entry.mark]{a}}, anchor=fn-1]
A note of the primary file.

\flmxItem[label={\flmxSpan[style=superscript, role=flm.xnotes.entry.mark]{b}}, anchor=fn-2]
A note of the included file.

\flmxItem[label={\flmxSpan[style=superscript, role=flm.xnotes.entry.mark]{c}}, anchor=fn-3]
A note of the primary file after the included file.
\end{flmxList}
\end{flmxBlock}

\end{document}
